\documentclass[a4paper,12pt]{article}
\usepackage[utf8]{inputenc}
\usepackage[T2A]{fontenc}
\usepackage[russian]{babel}

\begin{document}

\title{Генетические и меметические алгоритмы для решения задачи TSP}
\author{Пучков Пётр}
\date{\today}
\maketitle

\begin{abstract}
Настоящая работа посвящена применению генетических и меметических алгоритмов к NP-трудной задаче коммивояжёра (TSP). В ходе работы будут изучены различные виды мутаций, скрещивания и отбора. Кроме того будут рассмотрены гибридные методы, сочетающие в себе генетические алгоритмы и локальную оптимизацию. Изученные алгоритмы будут реализованы и протестированы. На основе полученных во время тестирования данных, будет проведён сравнительный анализ эффективности выбранных эвристик.

\end{abstract}

\pagebreak

\section{План работы}

\begin{enumerate}
    \item \textbf{Изучение основ.}
    \begin{itemize}
        \item Прочитаю рекомендованные главы~7--9 в Handbook of Metaheuristics, при необходимости -- дополнительные источники.
        \item Сделаю краткий обзор основных подходов к решению: точные, эвристические, метаэвристические.
    \end{itemize}

    \item \textbf{Освоение генетических и меметических алгоритмов.}
    \begin{itemize}
        \item Разберу базовые компоненты генетического алгоритма: представление решений, популяция, селекция, скрещивание, мутация.
        \item Пойму, как работает меметический алгоритм, т.е. как добавляется локальный поиск для улучшения каждого решения.
    \end{itemize}

    \item \textbf{Обзор и выбор операторов.}
    \begin{itemize}
        \item Изучу известные операторы селекции, скрещивания, мутации и локального поиска, применяемые для TSP.
    \end{itemize}

    \item \textbf{Программная реализация.}
    \begin{itemize}
        \item Реализую алгоритм вместе с архитектурой, позволяющую легко заменять эвристики.
    \end{itemize}

    \item \textbf{Экспериментальное сравнение.}
    \begin{itemize}
        % \item Заполучу тестовые данные и бейзлайн у В. Кондратюка
        \item Запущу алгоритмы на тестовых наборах данных, измерю качество решений, время работы, скорость сходимости.
        \item Построю графики и таблицы, проанализирую, как разные операторы и параметры влияют на результат.
    \end{itemize}
\end{enumerate}

\end{document}


